\documentclass[10pt,a4paper]{amsart}
\usepackage[margin=19mm]{geometry}
\usepackage{amsmath,amssymb,mathtools}
\usepackage[scr=rsfs,cal=euler]{mathalfa}
\usepackage{xfrac}
\usepackage[unicode,hidelinks,bookmarks=false]{hyperref}
\newcommand{\ie}{i.e.\ }
\newcommand{\cf}{cf.\ }
\newcommand{\resp}{respectively\ }
\newcommand{\Subsection}{\subsection}
\providecommand{\nobreakdash}{\nobreak\hbox{-}}
\setlength{\emergencystretch}{2em}
\multlinegap0pt
\usepackage{bm}
\usepackage{bbm}
\usepackage{dsfont}
\usepackage{xspace}
\usepackage{cancel}
\usepackage{mathtools}
\usepackage{amsthm}
\makeatletter
\@ifundefined{lemma}{\newtheorem{lemma}{Lemma}}{}
\@ifundefined{theorem}{\newtheorem{theorem}{Theorem}}{}
\makeatother
\newcommand{\Man}{\mathcal{M}}
\newcommand{\RM}{\mathrm{RM}}
\newcommand{\Vol}{\mathrm{Vol}}
\newcommand{\normt}[1]{%
	\vert\kern-0.9pt\vert\kern-0.9pt\vert #1
	\vert\kern-0.9pt\vert\kern-0.9pt\vert}
\title[Finite element methods for isometric embedding of Riemannian]{Finite element methods for isometric embedding of Riemannian manifolds}
\author{Guangwei Gao, Kaibo Hu, Buyang Li, Ganghui Zhang}
\date{Source: arXiv:2602.18722v1 (CC BY 4.0)}
\begin{document}
\maketitle

\noindent\emph{Source and licence.} Finite element methods for isometric embedding of Riemannian manifolds. Version arXiv:2602.18722v1 (CC BY 4.0), distributed under the Creative Commons Attribution 4.0 International licence, as recorded in the arXiv licence metadata for this exact version and shown on the abstract page https://arxiv.org/abs/2602.18722v1. Full rights notice: licenses/arxiv-2602.18722-CC-BY-4.0.txt.\par\smallskip

\noindent\textit{Editorial scope.} Counted unit: the Lemma whose source label is \texttt{lm:esti\_ev} in the article (source0.tex lines 1183--1188, source label \texttt{lm:esti\_ev}), delivered together with the article's own complete original proof of it (source0.tex lines 1189--1297, one \texttt{proof} environment of 109 lines). No mathematics is written, repaired or re-derived: statement and proof are the article's own text line for line. Required context retained uncounted: eq:isoemb (lines 328--330); eq:pt\_isoemb (lines 344--348); thm:wely-emb (lines 363--371); eq:g\_ell (lines 756--761); eq:gh\_app (lines 764--766); eq:num (lines 776--788); thm:err (lines 816--832); eq:graphnorm\_esti (lines 1055--1059); eq:inv (lines 1062--1064); eq:err-a (lines 1134--1142); eq:err-b (lines 1134--1142); eq:t* (lines 1153--1158); eq:t*:cor (lines 1161--1163). Every definition, display and label of those lines is retained unchanged. Omitted from this package and not redistributed: 1--327, 331--343, 349--362, 372--755, 762--763, 767--775, 789--815, 833--1054, 1060--1061, 1065--1133, 1143--1152, 1159--1160, 1164--1182, 1298--2101. Nothing inside a retained excerpt is omitted or altered: every retained body is delivered exactly as the article's lines are written, so no omission receipt of any kind accompanies this package. Citations: the retained excerpts carry no citation command at all, so no bibliography entry of the article is transcribed and no reference is redistributed. Reference adaptation: none was needed and none was made; every reference inside the delivered unit resolves inside it, so no unresolved reference or citation is produced. Printed slips kept literal: no printed word, symbol, hypothesis or number of the counted statement or of its proof was altered. Layout and class: the article's own class is \texttt{amsart} with packages including amsmath, amsthm, amsfonts, amssymb, bm, mathrsfs, color, xcolor, graphicx, comment, commath, mathtools, ulem, geometry; the wrapper is the lane's pinned \texttt{amsart} editorial wrapper at 10pt on A4, which restates the theorem-like environments the retained text needs and adds the article's own macro definitions that the retained text needs (\texttt{\textbackslash Man}, \texttt{\textbackslash RM}, \texttt{\textbackslash Vol}, \texttt{\textbackslash normt}), with extra wrapper packages (bm, bbm, dsfont, xspace, cancel, mathtools). The delivered numbering is the unit's own. Assets: no retained span contains a figure environment, a graphics-inclusion command, a PDF-page inclusion, a table or a TikZ picture; no image file is copied into this package, the package declares no asset dependency and no image file is redistributed with it. Licence: version arXiv:2602.18722v1 of this article is distributed under the Creative Commons Attribution 4.0 International licence, as recorded in the arXiv licence metadata for this exact version and shown on the abstract page https://arxiv.org/abs/2602.18722v1. A full rights notice is delivered at licenses/arxiv-2602.18722-CC-BY-4.0.txt. Characters: every retained excerpt is pure ASCII, so the delivered engine, XeLaTeX with its default Latin Modern text font, reports no missing character.
\begin{equation}\label{eq:isoemb}
\mathrm{d}r(t)\odot \mathrm{d}r(t)=g(t),\qquad t\in[0,T].
\end{equation}

\begin{equation}\label{eq:pt_isoemb}
2\,\mathrm{d}r \odot \mathrm{d}v = \partial_t g,
\quad \text{ where } \quad 
\partial_t r = v,
\end{equation}

\begin{theorem}\label{thm:wely-emb}
Let \(g(t)\) be a smooth Riemannian metric on \(\Man\), depending smoothly on \(t\in[0,T]\), with smooth time derivative \(\partial_t g(t)\).
Assume that the Gaussian curvature of \(g(t)\) remains strictly positive for all \(t\in[0,T]\).
Then there exists a {unique} smooth embedding flow \(r(t):\Man\to\mathbb{R}^3\), depending smoothly on \(t\in[0,T]\), such that \(r(t)\) satisfies~\eqref{eq:isoemb} and its time derivative \(\partial_t r\) satisfies~\eqref{eq:pt_isoemb} together with the orthogonality condition
\begin{equation}\label{eq:ptr_orth_RM}
\int_{\Man} \partial_t r \cdot (\alpha \times r + \beta)\,\Vol_{\Man} = 0,
\qquad \forall\, \alpha,\beta \in \mathbb{R}^3.
\end{equation}
\end{theorem}

\begin{equation}\label{eq:g_ell}
g^{(-\ell)} = a^*(\mathrm{d}r\odot \mathrm{d}r)
= \mathrm{d}r^{(-\ell)}\odot \mathrm{d}r^{(-\ell)} \in \Sigma,
\quad\text{and}\quad
\partial_t \, g^{(-\ell)} = 2\,\mathrm{d}r^{(-\ell)}\odot \mathrm{d}\partial_t r^{(-\ell)} \in \Sigma,
\end{equation}

\begin{equation}\label{eq:gh_app} 
    \|g_h(t) - g^{(-\ell)}(t)\|_{L^2(\Man_h)} \lesssim h^k, \qquad \|\partial_t g_h(t) - \partial_t g^{(-\ell)}(t)\|_{L^2(\Man_h)} \lesssim h^k. 
\end{equation}

\begin{subequations}\label{eq:num}  
\begin{align}
    \hspace{-28mm}
\label{eq:num-v}    
\partial_t r_h &= v_h, \\    
\label{eq:num-a} 
2(\mathrm{d} r_h \odot \mathrm{d} v_h, \mathrm{d} r_h \odot \mathrm{d} q_h)_{\Man_h} 
+ (\lambda_h, q_h)_{\Man_h}  &= 
(\partial_t g_h, \mathrm{d} r_h \odot \mathrm{d} q_h)_{\Man_h}, \\ 
\label{eq:num-b} 
(v_h, \mu_h)_{\Man_h} &= 0,
\end{align}
\end{subequations}

\begin{theorem}\label{thm:err}
Under the assumptions of Theorem~\ref{thm:wely-emb}, let \(r(t)\) be the smooth solution obtained therein.
Assume that the polynomial degree satisfies \(k=k_g\ge 5\).
Then there exists \(h_0>0\) such that, for all \(0<h\le h_0\), the numerical scheme~\eqref{eq:num} admits a unique solution
\(\bigl(r_h(t),\partial_t r_h(t),\lambda_h(t)\bigr)\) with \(\lambda_h(t)\equiv 0\), which exists uniquely on \([0,T]\) and satisfies
\[
r_h \in C^1([0,T];V_h^3).
\]
Moreover, the following error estimate holds:
\[
\normt{r_h^{(\ell)}(t)-r(t)}_{\Man}
=\Big(\|r_h^{(\ell)}(t)-r(t)\|_{L^2(\Man)}^2
+\|{\rm D}_r\bigl(r_h^{(\ell)}(t)-r(t)\bigr)\|_{L^2(\Man)}^2\Big)^{1/2}
\;\lesssim\; h^{k},
\]
for all \(t\in[0,T]\).
\end{theorem}

\begin{equation}\label{eq:graphnorm_esti}
\normt{v_h}_h \leq \normt{\Pi_{\RM} v_h}_h + \normt{ v_h - \Pi_{\RM} v_h}_h
\lesssim \|\Pi_{\RM} v_h \|_{L^2(\Man_h)} +  
\|\mathrm{d} r_h^* \odot \mathrm{d} v_h\|_{L^2(\Man_h)},
\end{equation}

\begin{equation}\label{eq:inv}
\|v_h\|_{W^{1,\infty}(\Man_h)} \lesssim h^{-1}\,\|v_h\|_{H^1(\Man_h)} \lesssim h^{-2}\,\|v_h\|_{L^2(\Man_h)} \leq h^{-2}\,\normt{v_h}_h, \quad \forall v_h \in V_h^3.
\end{equation}

\begin{align}
2(\partial_t(   \mathrm{d} r_h^* \odot\mathrm{d} e_r), \mathrm{d} r_h^* \odot \mathrm{d} q_h)_{\Man_h} 
 &= ( \partial_t g_h - 2\, \mathrm{d}r_h \odot \mathrm{d} \partial_t r_h,\ \mathrm{d}e_r \odot \mathrm{d}q_h )_{\Man_h}  \notag \\
\label{eq:err-a} 
& \quad
- ( \partial_t (\mathrm{d}e_r \odot \mathrm{d}e_r),\ \mathrm{d}r_h^* \odot \mathrm{d}q_h )_{\Man_h} - (d_v, q_h)_{\Man_h}, \\
\label{eq:err-b} 
( \partial_t e_r, \mu_h^* )_{\Man_h} &= - (\partial_t r_h, \mu_h^{(1)} \times e_r)_{\Man_h} - (d_{\lambda}, \mu_h^*)_{\Man_h},
\end{align}

\begin{equation}\label{eq:t*} 
\begin{split}
    t^*=&\sup\left\{ t \in [0, T]: \eqref{eq:num} \text{ has a solution on } [0,t],  \right. \\
     &\quad\quad \text{and } \normt{e_r(s)}_h \leq h^{4 + \varepsilon}, \text{ for } 0 \leq s \leq t\ \},
\end{split} 
\end{equation}

\begin{equation}\label{eq:t*:cor}
    \|  e_{r}(s) \|_{W^{1,\infty}(\Man_h)}\lesssim h^{2+\varepsilon}, \quad s \in [0,t^*].
\end{equation}

\begin{lemma}[Estimate of \( e_v \)]\label{lm:esti_ev}
    Under the assumptions of Theorem \ref{thm:err}, there exists a constant $h_0 > 0$ such that the following estimate holds uniformly for $0<h\leq h_0$ and $t\in[0,t^{*}]${\rm:}
    \[
    \normt{e_v(t) }_h\lesssim h^{-1} \normt{e_r(t)}_h + \normt{d_v(t)}_{h,*} + \|d_\lambda(t)\|_{L^2(\Man_h)} + h^{k+1}.
    \]
\end{lemma}

\begin{proof}
Decomposing 
$ \partial_t( \mathrm{d} r_h^* \odot\mathrm{d} e_r) =  \mathrm{d} r_h^* \odot\mathrm{d} e_v 
+  \mathrm{d} \partial_t r_h^* \odot\mathrm{d} e_r $ and choosing $q_h = e_{v}$ in \eqref{eq:err-a}, 
it holds that
\[ 
\begin{aligned}
    2\|\mathrm{d} e_{v} \odot \mathrm{d} r_h^{*}\|_{L^2(\Man_h)}^2 
&=(\partial_t g_h - 2 \mathrm{d} r_h \odot \mathrm{d} \partial_t r_h, \mathrm{d} e_{r} \odot \mathrm{d} e_{v} )_{\Man_h}
- 2 ( \mathrm{d} e_{r} \odot \mathrm{d} e_{v}, \mathrm{d} r_h^{*} \odot \mathrm{d} e_{v} )_{\Man_h} \\
&\quad  - ( d_v, e_{v} )_{\Man_h}
- 2 ( \mathrm{d} e_{r} \odot \mathrm{d} \partial_t r_h^{*}, \mathrm{d} r_h^{*} \odot \mathrm{d} e_{v})_{\Man_h}\\ 
& =: J_1 + J_2 + J_3 + J_4.
\end{aligned}
\]
For the term $J_1$, we first use \eqref{eq:g_ell} and \eqref{eq:gh_app}, together with the approximation property of $r_h^*$ to deduce that 
\begin{equation}\label{eq:pt_gh-drhdrh} 
\begin{aligned}
   & \| \partial_t g_h - 2 \mathrm{d} r_h \odot \mathrm{d} \partial_t r_h \|_{L^2(\Man_h)} \\ 
 \lesssim &\, \|\mathrm{d}r_h^*\odot\mathrm{d}\partial_t r_h^* - \mathrm{d}r_h\odot\mathrm{d}\partial_t r_h\|_{L^2(\Man_h)} + \| \partial_t g_h -  \partial_t g^{(-\ell)}  \|_{L^2(\Man_h)} \\ 
 & + \| \mathrm{d}r^{(-\ell)}\odot\mathrm{d}\partial_t r^{(-\ell)}
 - \mathrm{d} r_h^* \odot \mathrm{d} \partial_t r_h^*\|_{L^2(\Man_h)} \\ 
\lesssim &\,  \| \mathrm{d} r_h^{*} \odot \mathrm{d} e_{v}  \|_{L^2(\Man_h)}
+ \|\mathrm{d} e_{r} \odot \mathrm{d} e_{v} \|_{L^2(\Man_h)} 
+ \|\mathrm{d} e_{r} \odot \mathrm{d} \partial_t r_h^{*}\|_{L^2(\Man_h)} + h^k \\ 
\lesssim & 
\, \normt{ e_v }_{h} + \| e_r \|_{W^{1,\infty}(\Man_h)} \| e_v \|_{H^1(\Man_h)} 
+ \| e_r \|_{H^1(\Man_h)} + h^k\\
\lesssim & \, 
(1 + h^{-1}\|e_r\|_{W^{1, \infty}(\Man_h)})\normt{ e_v }_{h} + h^{-1} \normt{e_r}_h + h^k 
\quad\quad\quad\quad (\text{by inverse estimate \eqref{eq:inv}})
\\ 
\lesssim &
\, (1 + h^{1+\varepsilon})\normt{ e_v }_{h} + h^{-1} \normt{e_r}_h + h^k,
\qquad\qquad\qquad\qquad\qquad
 (\text{by estimate \eqref{eq:t*:cor}}). 
\end{aligned}
\end{equation} 
Substituting the estimate \eqref{eq:pt_gh-drhdrh}, we obtain the estimate for $J_1$: 
\[ \begin{aligned}
    J_1 &\le  
\|  \mathrm{d} e_{r} \odot \mathrm{d} e_{v} \|_{L^2(\Man_h)}\| \partial_t g_h - 2 \mathrm{d} r_h \odot \mathrm{d} \partial_t r_h \|_{L^2(\Man_h)}\\
&\lesssim \|  e_{r} \|_{W^{1,\infty}(\Man_h)} \| e_v \|_{H^1(\Man_h)}
(\normt{ e_v }_{h} + h^{-1} \normt{e_r}_h + h^k ) \\ 
&\lesssim h^{-1}\|  e_{r} \|_{W^{1,\infty}(\Man_h)} \normt{e_v}_h 
(\normt{ e_v }_{h} + h^{-1} \normt{e_r}_h + h^k)
\qquad \text{(by inverse estimate \eqref{eq:inv})}
\\ 
&\lesssim h^{1+\varepsilon} \normt{e_v}_h
(\normt{ e_v }_{h} + h^{-1} \normt{e_r}_h + h^k )
\qquad\qquad\qquad\qquad \text{(by estimate \eqref{eq:t*:cor})}\\
& \leq (Ch^{1+\varepsilon} + \delta ) \normt{e_v}_h^2 
+ C_{\delta}( h^{2\varepsilon} \normt{e_r}_h^2 + h^{2k+2+2\varepsilon} ),
\end{aligned}
\]
where $\delta>0$ can be chosen arbitrarily small, and $C_{\delta}>0$ is a constant (independent of $h$) depending only on $\delta$. 
Similarly, by the inverse estimate \eqref{eq:inv} and \eqref{eq:t*:cor}, we obtain
\[
J_2 \lesssim
\| e_r \|_{W^{1,\infty}(\Man_h)} \| e_v \|_{H^1(\Man_h)} \normt{e_v}_h
\lesssim h^{1+\varepsilon} \normt{e_v}_h^2.
\]
For the estimates of \(J_3\) and \(J_4\), we have
\begin{subequations}
\begin{align}
J_3 & \leq \delta \normt{e_v}_h^2 + C_{\delta} \normt{d_v}_{h,*}^2, \notag \\
J_4 &\lesssim \|e_r\|_{H^1(\Man_h)}\normt{e_v}_h
\leq \delta \normt{e_v}_h^2 +
C_{\delta}h^{-2}\normt{e_r}_h^2. \notag
\end{align}
\end{subequations}
Collecting the bounds for $J_1,\ldots,J_4$, we have
\begin{equation}\label{eq:devdr}
    \|\mathrm{d} e_{v} \odot \mathrm{d} r_h^{*}\|_{L^2(\Man_h)}^2
\leq 
C_{\delta}\big( 
h^{-2}\normt{e_r}_h^2 + \normt{ d_v }_{h,*}^2 + h^{2(k+1)}\big) 
+ C(\delta + h^{1+\varepsilon}) \normt{e_v}^2_h.
\end{equation}
Taking $\mu_h^* = \Pi_{\RM} e_v = e_{v, \RM}^{(1)} \times r_h^* +  e_{v, \RM}^{(2)}$ in \eqref{eq:err-b}, we obtain 
\[ 
    (e_v, \Pi_{\RM} e_v)_{\Man_h} 
     = - (\partial_t r_h^*,  e_{v, \RM}^{(1)} \times e_r)_{\Man_h} 
    - (e_v, e_{v, \RM}^{(1)} \times e_r) _{\Man_h}
    - (d_{\lambda}, \Pi_{\RM} e_v)_{\Man_h}. 
\]
Since $\Pi_{\RM}: V_h^3 \to \RM[r_h^*]$ is the $L^2$-orthogonal projection, we deduce that
\begin{equation}\label{eq:e_v,RM}
\begin{aligned}
        \|\Pi_{\RM} e_v \|_{L^2(\Man_h)}^2 & \lesssim 
        \| e_v \|_{L^2(\Man_h)} \| e_r \|_{L^2(\Man_h)} + 
        \| e_r \|_{L^\infty(\Man_h)} \| e_v \|_{L^2(\Man_h)}^2\\
        & \quad + \| e_v \|_{L^2(\Man_h)} \| d_{\lambda} \|_{L^2(\Man_h)} \\ 
        & \leq C_{\delta} (\normt{e_r}_h^2 + \| d_{\lambda} \|_{L^2(\Man_h)}^2) + C(\delta + h^{3+ \varepsilon}) \normt{e_v}_h^2,
\end{aligned}
\end{equation}
where we have used  \eqref{eq:t*} and  the inverse estimate  $ \| e_r \|_{L^\infty(\Man_h)} \lesssim h^{-1} \| e_r \|_{L^2(\Man_h)} \lesssim h^{3 + \varepsilon} $.
Combining \eqref{eq:graphnorm_esti},  \eqref{eq:devdr} and \eqref{eq:e_v,RM}, we have 
\[ \begin{aligned}
    \normt{e_v}^2_h &\lesssim 
    \|\Pi_{\RM} e_v\|_{L^2(\Man_h)}^2 +  \|\mathrm{d} e_{v} \odot \mathrm{d} r_h^{*}\|_{L^2(\Man_h)}^2 \\ 
    & \leq C_{\delta}\big( h^{-2}\normt{e_r}_h^2  + \| d_{\lambda} \|_{L^2(\Man_h)}^2 + \normt{d_v}_{h,*}^2 + h^{2(k+1)} \big)
 + C(\delta + h^{1+\varepsilon}) \normt{e_v}_h^2. 
\end{aligned} \]
Noting that, for sufficiently small $h$ and $\delta$, we can choose them so that
$C(\delta + h^{1+\varepsilon}) \leq \tfrac12$, the term
$C(\delta + h^{1+\varepsilon})\,\normt{e_v}_h^2$ can be absorbed into the left-hand side.
This completes the proof of Lemma \ref{lm:esti_ev}.
\end{proof}

\end{document}
